%% Modern LaTeX Academic Paper Template
%% Optimized for LuaLaTeX with Japanese support and modern bibliography

\documentclass[a4paper, 12pt, oneside]{article}

%% Language and font configuration
\usepackage{fontspec}
\usepackage{luatexja-fontspec}

% English fonts
\setmainfont{Latin Modern Roman}
\setsansfont{Latin Modern Sans}
\setmonofont{Latin Modern Mono}

% Japanese fonts (if needed)
\setmainjfont{Noto Serif CJK JP}
\setsansjfont{Noto Sans CJK JP}

%% Modern bibliography with BibLaTeX
\usepackage[
    backend=biber,
    style=authoryear-comp,
    sorting=nyt,
    url=false,
    doi=true,
    isbn=false,
    maxbibnames=3,
    maxcitenames=2,
    uniquelist=false,
    dashed=false
]{biblatex}
\addbibresource{references.bib}

%% Essential packages
\usepackage[utf8]{inputenc}
\usepackage[T1]{fontenc}
\usepackage{amsmath, amsfonts, amssymb}
\usepackage{graphicx}
\usepackage{booktabs}
\usepackage{microtype}
\usepackage{xcolor}

%% Page layout
\usepackage[
    a4paper,
    margin=2.5cm,
    top=3cm,
    bottom=3cm,
    includeheadfoot
]{geometry}

%% Headers and footers
\usepackage{fancyhdr}
\pagestyle{fancy}
\fancyhf{}
\fancyhead[L]{\textit{\@title}}
\fancyhead[R]{\textit{\@author}}
\fancyfoot[C]{\thepage}
\renewcommand{\headrulewidth}{0.4pt}

%% Hyperlinks and PDF metadata
\usepackage[
    colorlinks=true,
    linkcolor=blue!70!black,
    citecolor=red!70!black,
    urlcolor=blue!70!black,
    filecolor=magenta,
    bookmarks=true,
    bookmarksnumbered=true,
    pdfstartview=FitH,
    pdfpagemode=UseOutlines
]{hyperref}

%% Enhanced table support
\usepackage{array}
\usepackage{longtable}
\usepackage{multirow}

%% Code listings (if needed)
\usepackage{listings}
\lstset{
    basicstyle=\ttfamily\footnotesize,
    frame=single,
    numbers=left,
    numberstyle=\tiny,
    breaklines=true,
    showstringspaces=false,
    columns=flexible,
    commentstyle=\color{green!60!black},
    keywordstyle=\color{blue!80!black},
    stringstyle=\color{red!60!black}
}

%% Theorems and proofs
\usepackage{amsthm}
\newtheorem{theorem}{Theorem}[section]
\newtheorem{lemma}[theorem]{Lemma}
\newtheorem{proposition}[theorem]{Proposition}
\newtheorem{corollary}[theorem]{Corollary}
\theoremstyle{definition}
\newtheorem{definition}[theorem]{Definition}
\newtheorem{example}[theorem]{Example}
\theoremstyle{remark}
\newtheorem{remark}[theorem]{Remark}
\newtheorem{note}[theorem]{Note}

%% Enhanced figure and table captions
\usepackage[
    font=small,
    labelfont=bf,
    margin=1cm,
    format=hang
]{caption}

%% Appendix support
\usepackage[toc,page]{appendix}

%% Enhanced enumeration
\usepackage{enumitem}
\setlist[itemize]{noitemsep, topsep=0pt}
\setlist[enumerate]{noitemsep, topsep=0pt}

%% Custom commands for common notation
\newcommand{\R}{\mathbb{R}}
\newcommand{\N}{\mathbb{N}}
\newcommand{\Z}{\mathbb{Z}}
\newcommand{\Q}{\mathbb{Q}}
\newcommand{\C}{\mathbb{C}}
\newcommand{\E}{\mathbb{E}}
\newcommand{\Var}{\text{Var}}
\newcommand{\Cov}{\text{Cov}}

%% Document metadata (customize these)
\title{Your Paper Title Here}
\author{
    Your Name\thanks{Department, University, Email: your.email@university.edu} \\
    \textit{University Name} \\
    \and
    Co-Author Name\thanks{Different Department, Email: coauthor@university.edu} \\
    \textit{Co-Author Institution}
}
\date{\today}

%% PDF metadata
\hypersetup{
    pdftitle={Your Paper Title Here},
    pdfauthor={Your Name, Co-Author Name},
    pdfsubject={Academic Research Paper},
    pdfkeywords={keyword1, keyword2, keyword3, research},
    pdfcreator={LaTeX with LuaLaTeX}
}

%% Document begins
\begin{document}

%% Title page
\maketitle

%% Abstract
\begin{abstract}
This is the abstract of your paper. It should provide a concise summary of the research problem, methodology, key findings, and conclusions. The abstract should be self-contained and informative, typically ranging from 150-300 words depending on the journal requirements.

\smallskip
\noindent \textbf{Keywords:} keyword1; keyword2; keyword3; research methodology; academic writing
\end{abstract}

%% Main content
\section{Introduction}
\label{sec:introduction}

This is the introduction section. Here you should:
\begin{itemize}
    \item Provide background and motivation for your research
    \item Clearly state the research problem and objectives
    \item Review relevant literature (use citations like \textcite{sample2023} or \parencite{example2022})
    \item Outline the structure of the paper
\end{itemize}

Academic writing benefits from clear structure and proper citations. This template supports modern bibliography management with BibLaTeX and Biber.

\section{Literature Review}
\label{sec:literature}

In this section, you review existing work related to your research. Some key points to consider:

\begin{enumerate}
    \item Organize the literature thematically or chronologically
    \item Identify gaps in existing research
    \item Position your work within the broader academic context
\end{enumerate}

For example, \textcite{important2021} demonstrated that... while \textcite{another2020} argued for a different approach.

\section{Methodology}
\label{sec:methodology}

Describe your research methodology here. This might include:

\subsection{Data Collection}
Explain how you collected your data, including:
\begin{itemize}
    \item Sample size and selection criteria
    \item Data collection instruments
    \item Ethical considerations
\end{itemize}

\subsection{Analysis Methods}
Detail your analytical approach:

\begin{equation}
    \hat{\beta} = (X^T X)^{-1} X^T y
    \label{eq:ols}
\end{equation}

Where $\hat{\beta}$ represents the estimated coefficients, as shown in Equation~\ref{eq:ols}.

\section{Results}
\label{sec:results}

Present your findings here. Use tables and figures to illustrate key results.

\subsection{Descriptive Statistics}

Table~\ref{tab:descriptive} shows the descriptive statistics for our dataset.

\begin{table}[htbp]
    \centering
    \caption{Descriptive Statistics}
    \label{tab:descriptive}
    \begin{tabular}{lrrr}
        \toprule
        Variable & Mean & Std. Dev. & N \\
        \midrule
        Variable 1 & 25.4 & 5.2 & 150 \\
        Variable 2 & 18.9 & 3.8 & 150 \\
        Variable 3 & 42.1 & 8.7 & 150 \\
        \bottomrule
    \end{tabular}
    \footnotesize
    \textit{Note:} Sample descriptive statistics. Replace with your actual data.
\end{table}

\subsection{Main Findings}

Figure~\ref{fig:results} illustrates our main findings.

\begin{figure}[htbp]
    \centering
    \includegraphics[width=0.8\textwidth]{figures/placeholder.pdf}
    \caption{Main Results Visualization}
    \label{fig:results}
    \footnotesize
    \textit{Note:} Replace with your actual figure. The figure file should be placed in the \texttt{figures/} directory.
\end{figure}

\section{Discussion}
\label{sec:discussion}

Interpret your results in this section:
\begin{itemize}
    \item What do the results mean?
    \item How do they relate to existing literature?
    \item What are the implications?
    \item What are the limitations?
\end{itemize}

\section{Conclusion}
\label{sec:conclusion}

Summarize your contributions and suggest directions for future research.

\subsection{Key Contributions}
\begin{enumerate}
    \item First main contribution
    \item Second main contribution  
    \item Third main contribution
\end{enumerate}

\subsection{Future Work}
Outline potential extensions and improvements to your research.

%% Acknowledgments
\section*{Acknowledgments}
We thank [funding agencies, research assistants, colleagues, etc.] for their support and contributions to this research.

%% Bibliography
\printbibliography[heading=bibintoc, title=References]

%% Appendices (if needed)
\begin{appendices}
\section{Additional Results}
\label{app:additional}

Place supplementary material here if needed.

\section{Code Implementation}
\label{app:code}

Include relevant code snippets or refer to online repositories.

\end{appendices}

\end{document}